We will derive Theorem~\ref{thm:interiornonvanish} from  the following more technical statement.

\begin{theorem}\label{thm:localizedcase}
Let $\sigma\colon \Gamma \to 2^V$ be a quasi-geometric homology triangulation, and let $E$ be a face. Let $F \in \lk_{\Gamma}(E)$ be a face with $F \sqcup E$ interior. Suppose $\mathcal{A}_F = \emptyset$ and $F$ admits an interior partition $F = F_1 \sqcup F_2$. Assume that $F$ has no faces $G$ with $G \sqcup E$ interior and $ \vert G \vert < \vert F_1 \vert$. If $\vert F_1 \vert \le 2$, then $L(\Gamma, E)|_{F}$ is non-zero in degree $\vert F_1 \vert$. 
\end{theorem}

\begin{example}\label{example:nonrestriction}
The conclusion of Theorem~\ref{thm:localizedcase} can fail when $\vert F_1 \vert \ge 3$, even for $\mathcal{A}_F = E = \emptyset$.  Consider a geometric triangulation $\sigma \colon \Gamma \to 2^V$, where $V = \{v_1, \ldots, v_6 \}$ with a face $F = \{w_1, \ldots, w_6\}$ such that 
\[
\begin{array}{lll}
\sigma(w_1) = \{ v_1, v_3, v_6 \} & \sigma(w_2) = \{ v_1, v_4, v_5 \} & \sigma(w_3) = \{ v_2, v_3, v_5 \} \\
\sigma(w_4) = \{ v_2, v_4, v_6 \} & \sigma(w_5) = \{ v_3, v_4, v_5 \} & \sigma(w_6) = \{ v_3, v_5, v_6 \}
\end{array}
\]
Then $\mathcal{A}_F = \emptyset$, and $F$ admits an interior partition given by $F_1 = \{w_1, w_4, w_5\}$, $F_2 = \{w_2, w_3, w_6\}$. Then \eqref{e:restrict} gives generators and relations for $L(\Gamma, \emptyset)|_{F}$, and a linear algebra computation shows that $L(\Gamma, \emptyset)|_{F} = 0$. 
\end{example}
